% ============================================================
% CAPÍTULO 6 — Discusión
% ============================================================

\section{Evaluaci\'{o}n del desempe\~{n}o del modelo Seasonal Autoregressive Integrated Moving Average with Exogenous Variables (SARIMAX)}
\vspace{0.3cm}

\justify
El resultado m\'{a}s llamativo de la modelaci\'{o}n SARIMAX es que la b\'{u}squeda exhaustiva por AIC seleccion\'{o}, tanto para el cemento como para el ladrillo, el orden $(0,1,0)(0,0,0)_{12}$. Este orden equivale a un paseo aleatorio (\textit{random walk}) con variables ex\'{o}genas: la mejor predicci\'{o}n del precio del mes siguiente es, simplemente, el precio del mes actual m\'{a}s la contribuci\'{o}n lineal de las covariables. La ausencia de t\'{e}rminos autorregresivos ($p=0$), de media m\'{o}vil ($q=0$) y de componentes estacionales ($P=D=Q=0$) indica que, una vez diferenciada, la serie de precios se comporta como ruido blanco y que el modelo no encuentra estructura temporal adicional que pueda explotar. Este hallazgo es consistente con la din\'{a}mica de los mercados de insumos de construcci\'{o}n en Paraguay, donde los ajustes de precios son infrecuentes y tienden a mantenerse estables durante per\'{i}odos prolongados antes de registrar saltos discretos.

\justify
El desempe\~{n}o del SARIMAX sobre el conjunto de prueba difiri\'{o} significativamente entre materiales (Tablas \ref{tab:sarimax_metricas_prueba} y \ref{tab:sarimax_ladrillo_metricas_prueba}). Para el ladrillo, el RMSE de prueba fue de apenas Gs.~4,54, resultado que se explica porque durante los 22 meses del per\'{i}odo de evaluaci\'{o}n el precio se mantuvo constante en Gs.~650; en ese contexto, el paseo aleatorio ---que predice el \'{u}ltimo valor observado--- genera errores m\'{i}nimos y el $R^{2}$ resulta indefinido o cercano a cero por la escasa varianza del per\'{i}odo. Para el cemento, en cambio, el precio cay\'{o} de Gs.~60.000 a Gs.~55.000 en enero de 2024 y el modelo no anticip\'{o} ese ajuste, lo que se tradujo en un RMSE de prueba de Gs.~4.842,33 ---equivalente al tama\~{n}o del salto que el paseo aleatorio no pudo capturar---.

\justify
La validaci\'{o}n cruzada por ventana deslizante (\textit{rolling forecast}) ofrece una perspectiva m\'{a}s representativa del desempe\~{n}o real al abarcar todo el per\'{i}odo hist\'{o}rico, incluyendo los a\~{n}os de alta volatilidad asociados a la pandemia y a las variaciones hidrol\'{o}gicas. Bajo esta evaluaci\'{o}n, el $R^{2}$ sube a 0{,}7586 para el cemento (Tabla~\ref{tab:sarimax_metricas_cv}) y a 0{,}8193 para el ladrillo (Tabla~\ref{tab:sarimax_ladrillo_metricas_cv}), con RMSE de validaci\'{o}n cruzada de 1.920{,}78~Gs. y 9{,}52~Gs. respectivamente. Estos valores confirman la utilidad pr\'{a}ctica del SARIMAX para la estimaci\'{o}n presupuestaria a corto plazo cuando las series no presentan grandes saltos inesperados. El mejor ajuste del ladrillo puede atribuirse a que este material presenta per\'{i}odos prolongados de estabilidad, patr\'{o}n que el paseo aleatorio captura con naturalidad.

\justify
Un hallazgo cr\'{i}tico es que las variables ex\'{o}genas ---el nivel del r\'{i}o Paraguay y la variable binaria de cuarentena--- resultaron estad\'{i}sticamente no significativas en ambos modelos. Para el cemento, el coeficiente del nivel del r\'{i}o fue de $-55{,}34$~Gs./m con un $p$-valor de 0{,}895 (Tabla~\ref{tab:sarimax_params}); para el ladrillo, fue de $-1{,}93$~Gs./m con $p=0{,}302$ (Tabla~\ref{tab:sarimax_ladrillo_params}). La variable de cuarentena present\'{o} $p$-valores de 1{,}000 y 0{,}999, respectivamente. Esta no significancia no implica que dichas variables carezcan de influencia real sobre los precios, sino que el modelo lineal integrado ---al absorber la tendencia mediante la diferenciaci\'{o}n--- pierde la capacidad de detectar efectos que se manifiestan de forma no lineal, asim\'{e}trica o con retardos variables. El SARIMAX asume una relaci\'{o}n lineal y contempor\'{a}nea entre las covariables y la variable dependiente, supuesto que resulta insuficiente para capturar la complejidad de fen\'{o}menos como la pandemia de COVID-19, cuyo impacto sobre los precios se propag\'{o} a trav\'{e}s de m\'{u}ltiples canales (interrupciones log\'{i}sticas, ca\'{i}da de demanda, especulaci\'{o}n, reactivaci\'{o}n diferida) y con desfases temporales heterog\'{e}neos.

\justify
La insensibilidad del SARIMAX ante el escenario pand\'{e}mico lo confirma cuantitativamente: la diferencia entre el pron\'{o}stico con y sin la variable de cuarentena activa fue de apenas $-23{,}48$~Gs./mes para el cemento (Figuras~\ref{fig:sarimax_cemento_sin_covid} y \ref{fig:sarimax_cemento_con_covid}) y de $-0{,}82$~Gs./mes para el ladrillo (Figuras~\ref{fig:sarimax_ladrillo_sin_covid} y \ref{fig:sarimax_ladrillo_con_covid}), valores pr\'{a}cticamente despreciables frente a los niveles de precios de las series (Gs.~55.000 y Gs.~650, respectivamente). Esta limitaci\'{o}n constituye una de las motivaciones principales para explorar los modelos de aprendizaje profundo, que s\'{i} disponen de la flexibilidad no lineal necesaria para diferenciar escenarios.

\justify
En cuanto a los diagn\'{o}sticos estad\'{i}sticos, los residuos de ambos modelos SARIMAX superaron el test de Ljung-Box ($p>0{,}05$ en rezagos 10 y 20), lo que confirma la ausencia de autocorrelaci\'{o}n serial y valida que el modelo ha capturado adecuadamente la estructura din\'{a}mica lineal de las series. Sin embargo, el test de Shapiro-Wilk rechaz\'{o} la normalidad en ambos casos (Figuras~\ref{fig:sarimax_residuos} y \ref{fig:sarimax_ladrillo_residuos}), con valores de curtosis de 49{,}47 y 26{,}69 para cemento y ladrillo, respectivamente. Esta no normalidad ---atribuible a observaciones extremas durante el per\'{i}odo pand\'{e}mico--- no invalida las estimaciones puntuales, pero s\'{i} compromete la fiabilidad de los intervalos de confianza calculados bajo el supuesto gaussiano, que se ensanchan r\'{a}pidamente a medida que crece el horizonte de predicci\'{o}n. Para el cemento, el intervalo al 95\,\% pas\'{o} de Gs.~[49.535 -- 60.483] en el primer mes a Gs.~[28.089 -- 81.724] en el mes 24 (Tabla~\ref{tab:sarimax_pronostico}), un rango que cubre pr\'{a}cticamente todo el espectro hist\'{o}rico de la serie y limita la utilidad pr\'{a}ctica del pron\'{o}stico a largo plazo.

% ============================================================
\section{Evaluaci\'{o}n del desempe\~{n}o del modelo Long Short-Term Memory (LSTM)}
\vspace{0.3cm}

\justify
La primera aplicaci\'{o}n del modelo LSTM en este trabajo fue la predicci\'{o}n del nivel diario del r\'{i}o Paraguay, variable ex\'{o}gena central de los modelos de precios. Los resultados fueron notables: un RMSE de prueba de 0{,}0449~m ---menos de 5{,}3~cm de error sobre una serie que oscila entre $-2$ y $+8$~metros--- y un coeficiente de Nash-Sutcliffe (NSE) de 0{,}9992 (Tabla~\ref{tab:metricas_lstm_rio}). Este nivel de precisi\'{o}n, obtenido con una arquitectura de una sola capa LSTM de 512 unidades ocultas y una ventana de entrada de 180 d\'{i}as (Tabla~\ref{tab:hiperparametros_lstm_rio}), valida la capacidad de las redes LSTM para modelar series hidrol\'{o}gicas con din\'{a}micas estacionales complejas. La estrategia recursiva de predicci\'{o}n, en la que cada proyecci\'{o}n diaria se reincorpora como entrada para el paso siguiente, logr\'{o} generar un pron\'{o}stico a 730 d\'{i}as que reproduce correctamente los ciclos de crecida y estiaje del r\'{i}o (Figuras~\ref{fig:historico_prediccion} y \ref{fig:climatologia_prediccion}), aunque con una amplitud de oscilaci\'{o}n ligeramente atenuada respecto a la climatolog\'{i}a hist\'{o}rica ---comportamiento esperable en horizontes de proyecci\'{o}n extendidos donde la estrategia recursiva tiende a regresar las predicciones hacia la media---.

\justify
Esta predicci\'{o}n del nivel del r\'{i}o se utiliz\'{o} como variable ex\'{o}gena en todos los modelos de precios subsiguientes. Es importante se\~{n}alar que cualquier error en la proyecci\'{o}n hidrol\'{o}gica se propaga a las predicciones de precios, de modo que la cadena de modelaci\'{o}n introduce una fuente adicional de incertidumbre que no queda reflejada en los intervalos de confianza de los modelos de precios. No obstante, dado que el RMSE del modelo del r\'{i}o representa menos del 0{,}6\,\% del rango total de la variable, esta propagaci\'{o}n de errores puede considerarse marginal.

\justify
Para la predicci\'{o}n del precio del cemento, Optuna seleccion\'{o} una arquitectura compacta: una capa LSTM unidireccional con 64 unidades ocultas, sin \textit{dropout} ni planificador de tasa de aprendizaje, y con una ventana de entrada de 6 meses (Tabla~\ref{tab:lstm_cemento_hiper}). En el escenario sin pandemia, el modelo totaliz\'{o} 18.241 par\'{a}metros entrenables y alcanz\'{o} un RMSE de prueba de 5.432{,}05~Gs.\ (Tabla~\ref{tab:lstm_cemento_metricas}). En el escenario con pandemia, Optuna seleccion\'{o} una arquitectura bidireccional con \textit{dropout} y planificador ReduceLROnPlateau, resultando en 36.481 par\'{a}metros y un RMSE de prueba de 4.394{,}96~Gs. ---el mejor resultado individual para el cemento en este estudio---. Sin embargo, el an\'{a}lisis de residuos revel\'{o} autocorrelaci\'{o}n estad\'{i}sticamente significativa en los rezagos 10 ($p=0{,}017$) y 20 ($p=0{,}042$) del test de Ljung-Box (Figura~\ref{fig:lstm_cemento_acf_residuos}). Aunque los $p$-valores son cercanos al l\'{i}mite de significancia, este resultado indica que el modelo deja estructura temporal sin capturar en sus residuos, lo que podr\'{i}a traducirse en subestimaci\'{o}n de la incertidumbre cuando se emplean los intervalos de Monte Carlo Dropout. En el escenario sin pandemia, la ausencia de \textit{dropout} y de \textit{scheduler}, seleccionada por Optuna como configuraci\'{o}n \'{o}ptima, sugiere que la regularizaci\'{o}n L2 ($\lambda=10^{-5}$) fue suficiente para controlar el sobreajuste, pero no para eliminar toda la dependencia temporal en los errores.

\justify
El modelo LSTM del ladrillo constituye el mejor resultado individual de esta arquitectura. A diferencia del cemento, Optuna seleccion\'{o} una LSTM \textbf{bidireccional} con \textit{dropout} de 0{,}2, \textit{dropout} recurrente de 0{,}05 y planificador StepLR (Tabla~\ref{tab:lstm_ladrillo_hiper}), lo que result\'{o} en 36.481 par\'{a}metros entrenables ---aproximadamente el doble que el modelo del cemento sin pandemia---. El RMSE de prueba de 6{,}22~Gs.\ (escenario con pandemia) y 6{,}62~Gs.\ (sin pandemia) (Tabla~\ref{tab:lstm_ladrillo_metricas}) es notablemente bajo: representa un error relativo inferior al 2\,\% del rango hist\'{o}rico de la serie. El an\'{a}lisis de residuos mostr\'{o} un comportamiento m\'{a}s favorable que el del cemento: el test de Ljung-Box a 10 rezagos no rechaz\'{o} la hip\'{o}tesis nula ($p=0{,}237$), aunque s\'{i} lo hizo a 20 rezagos ($p<0{,}001$), indicando la presencia de patrones de dependencia en rezagos distantes que podr\'{i}an estar asociados a la estacionalidad no capturada completamente por las variables de entrada.

\justify
La selecci\'{o}n de una arquitectura bidireccional para el ladrillo es un hallazgo relevante. La bidireccionalidad permite al modelo procesar la secuencia temporal en ambas direcciones dentro de la ventana de entrada, lo que duplica la cantidad de informaci\'{o}n contextual disponible en cada paso. El hecho de que Optuna haya preferido esta configuraci\'{o}n para el ladrillo ---pero no para el cemento--- sugiere que la serie de precios del ladrillo, con sus ajustes escalonados y per\'{i}odos prolongados de estabilidad, se beneficia de la capacidad de la red bidireccional para ``mirar hacia adelante'' dentro de la ventana de contexto y anticipar los momentos de cambio. En contraste, la tendencia m\'{a}s continua del cemento se modela adecuadamente con una arquitectura unidireccional m\'{a}s simple.

\justify
En cuanto a la diferenciaci\'{o}n de escenarios pand\'{e}micos, el modelo LSTM mostr\'{o} un comportamiento consistente para ambos materiales: bajo el escenario con cuarentena activa, el l\'{i}mite superior del rango de predicci\'{o}n del cemento subi\'{o} de 62.197~Gs.\ a 73.374~Gs.\ ---un incremento del 18{,}1\,\%--- (Figuras~\ref{fig:lstm_cemento_pred_sin_covid} y \ref{fig:lstm_cemento_pred_con_covid}), mientras que para el ladrillo pas\'{o} de 706~Gs.\ a 757~Gs.\ ---un incremento del 7{,}2\,\%--- (Figuras~\ref{fig:lstm_ladrillo_pred_sin_covid} y \ref{fig:lstm_ladrillo_pred_con_covid}). Ambas direcciones son consistentes con lo ocurrido durante la pandemia real de 2020--2022, cuando las disrupciones en la cadena de suministro provocaron aumentos sostenidos en los precios de materiales de construcci\'{o}n. La menor magnitud del efecto en el ladrillo puede atribuirse a que este material, de producci\'{o}n predominantemente local con cadenas log\'{i}sticas m\'{a}s cortas, es menos sensible a las interrupciones del comercio internacional que afectaron particularmente al cemento.

% ============================================================
\section{Evaluaci\'{o}n del desempe\~{n}o del modelo Gated Recurrent Unit (GRU)}
\vspace{0.3cm}

\justify
La arquitectura GRU, que emplea dos compuertas (actualizaci\'{o}n y reinicio) frente a las tres de la LSTM (entrada, olvido y salida), fue evaluada como alternativa con menor complejidad param\'{e}trica. Para el cemento, Optuna seleccion\'{o} una configuraci\'{o}n con una capa GRU unidireccional de 64 unidades ocultas, \textit{dropout} de 0{,}1, optimizador Adam con $\eta=0{,}00732$ y planificador ReduceLROnPlateau (Tabla~\ref{tab:gru_cemento_arq}), resultando en 13.889 par\'{a}metros entrenables. El RMSE de prueba fue de 4.964{,}27~Gs.\ (Tabla~\ref{tab:gru_cemento_metricas}), superior al obtenido por la LSTM en el escenario con pandemia (4.394{,}96~Gs.). La selecci\'{o}n de \textit{dropout} por parte de Optuna ---ausente en el modelo LSTM del mismo material--- sugiere que la GRU, pese a su menor n\'{u}mero de par\'{a}metros, present\'{o} indicios de sobreajuste durante la b\'{u}squeda, posiblemente porque la simplificaci\'{o}n de compuertas reduce la capacidad de la red para gestionar la memoria de largo plazo en la serie del cemento, que exhibe una tendencia continua.

\justify
En cuanto al efecto pand\'{e}mico, la GRU del cemento mostr\'{o} una respuesta consistente con la LSTM: bajo el escenario con cuarentena, el l\'{i}mite superior del rango de predicci\'{o}n subi\'{o} de Gs.~63.852 a Gs.~65.337 ---un incremento del 2{,}3\,\%--- (Figuras~\ref{fig:gru_cemento_pred_sin_covid} y \ref{fig:gru_cemento_pred_con_covid}). Aunque la magnitud del efecto estimado por la GRU es considerablemente menor que el 18{,}1\,\% proyectado por la LSTM, ambos modelos coinciden en la direcci\'{o}n del impacto pand\'{e}mico sobre el cemento: una presi\'{o}n alcista coherente con las disrupciones en la cadena de suministros que elevaron los costos log\'{i}sticos durante 2020--2022.

\justify
El modelo GRU para el ladrillo, con 13.889 par\'{a}metros entrenables (Tabla~\ref{tab:gru_ladrillo_arq}), alcanz\'{o} un RMSE de prueba de 11{,}88~Gs.\ (Tabla~\ref{tab:gru_ladrillo_metricas}), resultado inferior tanto a la LSTM bidireccional del mismo material (6{,}22~Gs., 36.481 par\'{a}metros) como al modelo SARIMAX (4{,}54~Gs.). Este desempe\~{n}o indica que, para la serie del ladrillo ---caracterizada por ajustes escalonados y per\'{i}odos prolongados de estabilidad---, la arquitectura GRU unidireccional no logr\'{o} capturar los patrones de cambio con la misma efectividad que las alternativas evaluadas.

\justify
El an\'{a}lisis de residuos del GRU del ladrillo fue tambi\'{e}n el m\'{a}s favorable de todos los modelos de aprendizaje profundo evaluados: el test de Ljung-Box no rechaz\'{o} la hip\'{o}tesis nula de ausencia de autocorrelaci\'{o}n serial ni a 10 rezagos ($p=0{,}442$) ni a 20 rezagos ($p=0{,}081$) (Figura~\ref{fig:gru_ladrillo_acf_residuos}). Este resultado indica que los residuos del modelo se comportan esencialmente como ruido blanco, lo que valida la adecuaci\'{o}n de los intervalos de confianza estimados por Monte Carlo Dropout y otorga mayor fiabilidad a las proyecciones generadas.

\justify
La selecci\'{o}n de RMSprop como optimizador para el ladrillo, frente al Adam elegido en todos los dem\'{a}s modelos, es otro aspecto relevante. RMSprop, al adaptar la tasa de aprendizaje seg\'{u}n la magnitud del gradiente de cada par\'{a}metro, resulta particularmente efectivo para superficies de p\'{e}rdida con gradientes de magnitudes heterog\'{e}neas, como las que pueden generarse en series con saltos escalonados. Combinado con el planificador CosineAnnealingLR, que reduce la tasa de aprendizaje siguiendo un perfil cosinusoidal, el modelo logra un ajuste fino que explica, al menos en parte, la calidad superior de sus residuos.

\justify
En cuanto al comportamiento bajo escenarios pand\'{e}micos, el GRU del ladrillo proyect\'{o} un incremento de precios bajo cuarentena: el l\'{i}mite superior del rango subi\'{o} de Gs.~726 a Gs.~761 ---un aumento del 4{,}8\,\%--- (Figuras~\ref{fig:gru_ladrillo_pred_sin_covid} y \ref{fig:gru_ladrillo_pred_con_covid}), en l\'{i}nea con la interpretaci\'{o}n alcista de la LSTM para el mismo material. Esta coincidencia entre arquitecturas diferentes para el ladrillo refuerza la hip\'{o}tesis de que el efecto pand\'{e}mico sobre este material fue predominantemente inflacionario.

% ============================================================
\section{Comparaci\'{o}n entre modelos}
\vspace{0.3cm}

\justify
La Tabla~\ref{tab:comparacion_modelos} presenta una s\'{i}ntesis comparativa de los siete modelos evaluados en este trabajo, incluyendo las m\'{e}tricas de error, la complejidad param\'{e}trica y el estado de los diagn\'{o}sticos residuales.

\begin{table}[ht]
    \centering
    \caption{Comparaci\'{o}n integral de los modelos evaluados para la predicci\'{o}n de precios de materiales de construcci\'{o}n.}
    \label{tab:comparacion_modelos}
    \begin{tabular}{llrrll}
        \toprule
        \textbf{Modelo} & \textbf{Material} & \textbf{RMSE Test} & \textbf{Par\'{a}metros} & \textbf{Ljung-Box} & \textbf{Observaci\'{o}n} \\
        \midrule
        SARIMAX & Cemento  & 4.842{,}33~Gs.$^{a}$  & 3      & OK ($p>0{,}05$)               & $R^{2}_{\text{CV}}=0{,}76$ \\
        LSTM    & Cemento  & 4.394{,}96~Gs.$^{c}$  & 36.481 & Borderline ($p\approx0{,}02$) & Con pandemia; bidireccional \\
        GRU     & Cemento  & 4.964{,}27~Gs.        & 13.889 & OK ($p>0{,}05$)               & \textit{Dropout}=0{,}1 \\
        \midrule
        SARIMAX & Ladrillo & 4{,}54~Gs.$^{a}$      & 3      & OK ($p>0{,}05$)               & $R^{2}_{\text{CV}}=0{,}82$ \\
        LSTM    & Ladrillo & 6{,}22~Gs.$^{c}$      & 36.481 & Mixto$^{b}$                   & Con pandemia; bidireccional \\
        GRU     & Ladrillo & 11{,}88~Gs.           & 13.889 & OK ($p>0{,}05$)               & Unidireccional \\
        \midrule
        LSTM    & Nivel r\'{i}o & 0{,}0449~m        & 1.079.809 & ---             & NSE=0{,}9992 \\
        \bottomrule
        \multicolumn{6}{l}{\footnotesize $^{a}$ RMSE sobre conjunto de prueba; RMSE de CV: 1.920{,}78~Gs.\ (cemento) y 9{,}52~Gs.\ (ladrillo).} \\
        \multicolumn{6}{l}{\footnotesize $^{b}$ OK a 10 rezagos ($p=0{,}237$), significativo a 20 rezagos ($p<0{,}001$).} \\
        \multicolumn{6}{l}{\footnotesize $^{c}$ Mejor escenario (con pandemia); RMSE sin pandemia: 5.432{,}05~Gs.\ (cemento) y 6{,}62~Gs.\ (ladrillo).}
    \end{tabular}
\end{table}

\justify
La comparaci\'{o}n entre modelos revela varios hallazgos clave que merecen un an\'{a}lisis detallado.

\justify
\textbf{Para el cemento}, la LSTM con pandemia obtuvo el menor RMSE de prueba (4.394{,}96~Gs.), seguida por el SARIMAX en prueba directa (4.842{,}33~Gs.) y la GRU (4.964{,}27~Gs.). La m\'{e}trica m\'{a}s robusta del SARIMAX sigue siendo el RMSE de validaci\'{o}n cruzada (1.920{,}78~Gs.), que abarca el per\'{i}odo hist\'{o}rico completo con mayor variabilidad. Comparando directamente los RMSE de prueba, la LSTM con pandemia supera al SARIMAX en un 9{,}2\,\% de reducci\'{o}n de error ---una ganancia moderada que indica que ambos modelos son competitivos para la predicci\'{o}n del precio del cemento a corto plazo---.

\justify
\textbf{Para el ladrillo}, el SARIMAX fue el modelo m\'{a}s preciso en prueba directa (4{,}54~Gs.), seguido por la LSTM con pandemia (6{,}22~Gs.) y la GRU (11{,}88~Gs.). La superioridad del SARIMAX para este material se atribuye a la din\'{a}mica escalonada del ladrillo, cuyo precio se mantuvo constante en Gs.~650 durante todo el per\'{i}odo de evaluaci\'{o}n: en ese contexto, el paseo aleatorio resulta especialmente preciso. El desempe\~{n}o inferior de la GRU respecto a la LSTM para el ladrillo es consistente con la hip\'{o}tesis de que la arquitectura bidireccional de la LSTM ---que duplica la informaci\'{o}n contextual disponible--- es m\'{a}s efectiva para capturar los patrones de ajuste escalonado de esta serie.

\justify
\textbf{Ning\'{u}n modelo \'{u}nico domina a trav\'{e}s de ambos materiales.} Para el cemento, la LSTM con pandemia ofrece el mejor RMSE de prueba (4.394{,}96~Gs.), con el SARIMAX en validaci\'{o}n cruzada como referencia de largo plazo. Para el ladrillo, el modelo SARIMAX result\'{o} el m\'{a}s preciso en prueba directa (4{,}54~Gs.), seguido de la LSTM (6{,}22~Gs.) y la GRU (11{,}88~Gs.). Esta falta de dominancia universal es consistente con el principio del \textit{No Free Lunch} en aprendizaje autom\'{a}tico: no existe un algoritmo que sea \'{o}ptimo para todos los problemas, y la selecci\'{o}n del modelo debe estar guiada por las caracter\'{i}sticas espec\'{i}ficas de la serie a predecir.

\justify
La relaci\'{o}n entre complejidad param\'{e}trica y desempe\~{n}o muestra un patr\'{o}n mixto en este estudio. Para el cemento, la LSTM con pandemia (36.481 par\'{a}metros, bidireccional) fue el modelo m\'{a}s preciso, superando a la GRU (13.889 par\'{a}metros) y a la LSTM sin pandemia (18.241 par\'{a}metros). Para el ladrillo, en cambio, la LSTM bidireccional (36.481 par\'{a}metros, RMSE 6{,}22~Gs.) super\'{o} a la GRU unidireccional (13.889 par\'{a}metros, RMSE 11{,}88~Gs.), y el modelo cl\'{a}sico SARIMAX (3 par\'{a}metros, RMSE 4{,}54~Gs.) fue el mejor de todos. Este resultado ilustra que, para series con per\'{i}odos prolongados de estabilidad como la del ladrillo, la capacidad adicional de las redes neuronales no se traduce necesariamente en mayor precisi\'{o}n, y los modelos estad\'{i}sticos tradicionales pueden ser igualmente competitivos. La b\'{u}squeda autom\'{a}tica de Optuna demostr\'{o} su valor al seleccionar configuraciones que se adaptan a las caracter\'{i}sticas de cada serie.

\justify
Un aspecto diferenciador fundamental entre los modelos es su capacidad para capturar el efecto de la variable de cuarentena. Mientras que el SARIMAX se mostr\'{o} completamente insensible a esta variable (diferencias inferiores a Gs.~25 para ambos materiales), los modelos de aprendizaje profundo ---tanto LSTM como GRU--- generaron diferencias sustanciales entre los escenarios con y sin pandemia. Para el cemento, la LSTM proyect\'{o} un incremento del 18{,}1\,\% en el l\'{i}mite superior del rango, y para el ladrillo, la LSTM proyect\'{o} un incremento del 7{,}2\,\%. Esta capacidad de diferenciar escenarios constituye un valor a\~{n}adido pr\'{a}ctico significativo de los modelos de aprendizaje profundo: no solo predicen el nivel esperado del precio, sino que permiten evaluar la sensibilidad de las proyecciones ante perturbaciones ex\'{o}genas, funcionalidad esencial para la gesti\'{o}n de riesgo en presupuestos de obra.

% ============================================================
\section{Impacto de la variable de cuarentena por COVID-19}
\vspace{0.3cm}

\justify
El an\'{a}lisis del impacto de la variable binaria de cuarentena por COVID-19 sobre las proyecciones de precios constituye uno de los aportes centrales de este trabajo, dado que permite evaluar la capacidad de cada familia de modelos para capturar los efectos de perturbaciones ex\'{o}genas no lineales. Los resultados revelaron tres patrones diferenciados que merecen una discusi\'{o}n detallada.

\justify
\textbf{El modelo SARIMAX fue esencialmente insensible a la variable pand\'{e}mica.} Para el cemento, la diferencia entre los escenarios con y sin cuarentena fue de apenas $-23{,}48$~Gs./mes; para el ladrillo, de $-0{,}82$~Gs./mes. Estas cifras, inferiores al 0{,}05\,\% del nivel de precios de las series, confirman que el marco lineal integrado carece de la flexibilidad necesaria para representar el efecto estructural de la pandemia. La explicaci\'{o}n reside en la diferenciaci\'{o}n de la serie ($d=1$): al transformar la variable dependiente en cambios peri\'{o}dicos, el modelo absorbe los desplazamientos de nivel ---incluido el provocado por la pandemia--- y reduce la variable ex\'{o}gena a una contribuci\'{o}n marginal lineal que resulta estad\'{i}sticamente indistinguible de cero. Este resultado no implica que la pandemia no afect\'{o} los precios, sino que el efecto no puede ser capturado por un modelo que asume relaciones lineales y contempor\'{a}neas entre las variables.

\justify
\textbf{La LSTM proyect\'{o} de manera consistente un incremento de precios bajo el escenario pand\'{e}mico}, tanto para el cemento (+18{,}1\,\% en el l\'{i}mite superior) como para el ladrillo (+7{,}2\,\%). Esta interpretaci\'{o}n alcista es coherente con la din\'{a}mica real observada durante 2020--2022, cuando las disrupciones en las cadenas de suministro ---escasez de materias primas, aumento de costos de transporte, restricciones de importaci\'{o}n--- provocaron incrementos sostenidos en los precios de los materiales de construcci\'{o}n a nivel global. La LSTM, gracias a su mecanismo de compuertas que permite retener informaci\'{o}n relevante a lo largo de m\'{u}ltiples pasos temporales, aprendi\'{o} a asociar la se\~{n}al de cuarentena con la trayectoria ascendente que los precios siguieron durante ese per\'{i}odo. La menor magnitud del efecto en el ladrillo (7{,}2\,\% frente a 18{,}1\,\% del cemento) es consistente con la hip\'{o}tesis de que los materiales de producci\'{o}n local, con cadenas log\'{i}sticas m\'{a}s cortas, resultaron menos afectados por las interrupciones del comercio internacional.

\justify
\textbf{La GRU mostr\'{o} una respuesta alcista consistente para ambos materiales bajo el escenario pand\'{e}mico}: para el cemento, el l\'{i}mite superior del pron\'{o}stico ascendi\'{o} de Gs.~63.852 a Gs.~65.337 (un incremento del 2{,}3\,\%); para el ladrillo, de Gs.~726 a Gs.~761 (un incremento del 4{,}8\,\%). Aunque la magnitud de los incrementos estimados por la GRU es considerablemente menor que la proyectada por la LSTM, la direcci\'{o}n es coincidente en ambos modelos y ambos materiales: todos los modelos de aprendizaje profundo asociaron la variable de cuarentena con una presi\'{o}n alcista sobre los precios, coherente con el efecto dominante de disrupci\'{o}n de la cadena de suministros durante 2020--2022.

\justify
La consistencia en la direcci\'{o}n del efecto pand\'{e}mico entre LSTM y GRU para ambos materiales otorga mayor credibilidad a la hip\'{o}tesis alcista: la variable de cuarentena fue procesada por ambas arquitecturas como una se\~{n}al de presi\'{o}n ascendente sobre los precios, independientemente de los detalles de implementaci\'{o}n. Esta coincidencia refuerza la validez interna de los resultados y sugiere que el efecto estimado no es un artefacto de una arquitectura particular, sino una respuesta robusta del aprendizaje profundo ante el patr\'{o}n pand\'{e}mico en los datos. Desde el punto de vista metodol\'{o}gico, la representaci\'{o}n binaria del \textit{shock} pand\'{e}mico sigue siendo una simplificaci\'{o}n: una sola variable dicot\'{o}mica no puede capturar simult\'{a}neamente los distintos canales de transmisi\'{o}n ---disrupci\'{o}n de oferta, contracci\'{o}n de demanda, especulaci\'{o}n, reactivaci\'{o}n diferida--- que caracterizaron la pandemia de COVID-19. Un modelo m\'{a}s completo requerir\'{i}a variables que representen por separado cada uno de estos canales.

\justify
En s\'{i}ntesis, los modelos de aprendizaje profundo demostraron una capacidad significativamente superior al SARIMAX para diferenciar escenarios pand\'{e}micos, validando la inclusi\'{o}n de la variable de cuarentena como predictor relevante en el contexto del aprendizaje no lineal. La convergencia de LSTM y GRU hacia una direcci\'{o}n alcista para ambos materiales fortalece la robustez de esta conclusi\'{o}n, aunque la variable binaria resulta insuficiente para capturar la totalidad de los mecanismos de transmisi\'{o}n del \textit{shock} pand\'{e}mico.

% ============================================================
\section{Limitaciones del estudio}
\vspace{0.3cm}

\justify
A pesar de los resultados favorables obtenidos, este trabajo presenta limitaciones que deben considerarse al interpretar los hallazgos y que abren oportunidades para investigaciones futuras.

\justify
\textbf{Tama\~{n}o reducido de las series temporales.} Con 139 observaciones mensuales para las series de precios (enero 2014 -- julio 2025), el \textit{dataset} se encuentra en el l\'{i}mite inferior de lo que los modelos de aprendizaje profundo requieren para alcanzar su m\'{a}ximo potencial. Las redes neuronales recurrentes, dise\~{n}adas para aprender patrones a partir de grandes vol\'{u}menes de datos secuenciales, se ven forzadas en este contexto a operar con un n\'{u}mero de observaciones de entrenamiento (95 meses) apenas superior al n\'{u}mero de par\'{a}metros de los modelos m\'{a}s complejos. Esta restricci\'{o}n limita la profundidad de las arquitecturas que Optuna puede explorar y favorece configuraciones compactas ---como lo evidencia la selecci\'{o}n de modelos con una sola capa en todos los casos---. Con series m\'{a}s largas, ser\'{i}a posible evaluar arquitecturas m\'{a}s profundas, ventanas de entrada m\'{a}s amplias y estrategias de entrenamiento m\'{a}s sofisticadas que podr\'{i}an mejorar la capacidad predictiva.

\justify
\textbf{Representaci\'{o}n simplificada del efecto pand\'{e}mico.} La variable binaria de cuarentena, que toma valor 1 durante los meses de restricciones y 0 en el resto, constituye una simplificaci\'{o}n considerable de un fen\'{o}meno multifac\'{e}tico. La pandemia de COVID-19 afect\'{o} los precios de materiales de construcci\'{o}n a trav\'{e}s de m\'{u}ltiples canales simult\'{a}neos: contracci\'{o}n de la demanda por paralizaci\'{o}n de obras, disrupci\'{o}n de cadenas de suministro internacionales, escasez de insumos importados, aumento de costos de transporte y especulaci\'{o}n de mercado. Una variable binaria no puede diferenciar entre estos mecanismos, lo que explica las interpretaciones divergentes entre modelos (la LSTM capturando el efecto oferta y la GRU capturando, en algunos casos, el efecto demanda). Futuras investigaciones podr\'{i}an incorporar variables continuas que representen la intensidad de las restricciones, \'{i}ndices de movilidad, vol\'{u}menes de importaci\'{o}n o indicadores de actividad econ\'{o}mica sectorial.

\justify
\textbf{Propagaci\'{o}n de errores en la cadena de modelaci\'{o}n.} La predicci\'{o}n del nivel del r\'{i}o Paraguay, generada por el modelo LSTM univariado, se utiliz\'{o} como variable ex\'{o}gena en todos los modelos de precios. Aunque el error de esta predicci\'{o}n es peque\~{n}o (RMSE de 0{,}0449~m, menos del 0{,}6\,\% del rango de la variable), la estrategia recursiva a 730 d\'{i}as acumula incertidumbre a lo largo del horizonte de proyecci\'{o}n, y esta incertidumbre no queda reflejada en los intervalos de confianza de los modelos de precios. Un enfoque m\'{a}s riguroso requerir\'{i}a propagar la distribuci\'{o}n de errores del modelo hidrol\'{o}gico a trav\'{e}s de los modelos de precios, por ejemplo, mediante t\'{e}cnicas de simulaci\'{o}n que muestreen m\'{u}ltiples trayectorias posibles del nivel del r\'{i}o y generen un abanico de predicciones de precios condicionadas a cada una de ellas.

\justify
\textbf{Heterogeneidad en la comparaci\'{o}n de escenarios.} Los escenarios ``sin COVID'' y ``con COVID'' involucraron, en algunos casos, optimizaciones de Optuna independientes para cada condici\'{o}n, lo que implica que las arquitecturas e hiperpar\'{a}metros seleccionados pueden diferir entre ambos escenarios. Esta heterogeneidad dificulta una atribuci\'{o}n limpia del efecto pand\'{e}mico, dado que las diferencias entre pron\'{o}sticos pueden deberse tanto a la variable de cuarentena como a la configuraci\'{o}n diferente del modelo. Una comparaci\'{o}n m\'{a}s controlada requerir\'{i}a fijar la arquitectura y los hiperpar\'{a}metros, variando \'{u}nicamente el valor de la variable de cuarentena en el horizonte de predicci\'{o}n.

\justify
\textbf{Limitaciones de los intervalos de confianza.} Los intervalos de confianza de los modelos de aprendizaje profundo fueron estimados mediante Monte Carlo Dropout con $N=100$ simulaciones. Este m\'{e}todo, aunque pr\'{a}ctico y computacionalmente eficiente, asume que la variabilidad inducida por la desactivaci\'{o}n aleatoria de neuronas es un buen proxy de la incertidumbre epist\'{e}mica del modelo. Sin embargo, cuando los residuos presentan autocorrelaci\'{o}n ---como ocurre con la LSTM del cemento ($p=0{,}017$ a 10 rezagos) y con la LSTM del ladrillo ($p<0{,}001$ a 20 rezagos)---, los intervalos de Monte Carlo Dropout pueden subestimar la incertidumbre real, al no capturar la correlaci\'{o}n temporal de los errores. Para los modelos con residuos limpios, como la GRU del ladrillo, esta limitaci\'{o}n es menos relevante.

\justify
\textbf{Alcance limitado del SARIMAX.} El modelo SARIMAX alcanz\'{o} valores de RMSE razonables en validaci\'{o}n cruzada, lo que lo posiciona como una herramienta pr\'{a}ctica para la estimaci\'{o}n presupuestaria a corto plazo. Sin embargo, su incapacidad para capturar efectos no lineales y su insensibilidad ante la variable pand\'{e}mica limitan su utilidad en contextos donde se requiere modelar el impacto de perturbaciones estructurales. En mercados estables, el SARIMAX puede ser suficiente; en per\'{i}odos de crisis, los modelos de aprendizaje profundo ofrecen una ventaja cualitativa que trasciende las m\'{e}tricas de error puntuales.
